\section{Notation and glossary}\label{app:glossary}
Use this appendix whenever a symbol interrupts your reading. An index identifies
a particular object; it does not mean multiplication. Hats distinguish operators
from ordinary numbers. Bold symbols collect several numbers into a vector.

\renewcommand{\arraystretch}{1.12}
\begin{longtable}{@{}p{0.23\textwidth}>{\raggedright\arraybackslash}p{0.71\textwidth}@{}}
\toprule
Symbol or term & Meaning in this project\\
\midrule
\endfirsthead
\toprule
Symbol or term & Meaning in this project\\
\midrule
\endhead
\bottomrule
\endfoot
$x_i$ & A binary decision: select booking $i$ when one; leave it unselected when zero.\\
$\boldsymbol{x}$ & Logical vector $(x_0,\ldots,x_{n-1})$ describing a complete proposed schedule.\\
$q_i$ & Qubit carrying booking bit $i$ when measured.\\
$n$ & Number of booking variables/qubits in the direct encoding.\\
$\mathcal V$, $G_v$ & Vehicle set and booking-index set for vehicle $v$.\\
$\mathcal C$ & Unordered pairs of different-car bookings sharing a station and slot.\\
$c_i$, $\lambda$ & Booking score and waiting weight, with default $\lambda=0.25$ points/min.\\
$P$, $U$ & Constraint penalty and bound on any feasible schedule's operating score.\\
$C_0,L_i,Q_{ij}$ & Constant, linear, and upper-triangular pair coefficients in the expanded QUBO.\\
$z_i$ & Numerical sign $(-1)^{x_i}$: $+1$ for bit zero, $-1$ for bit one.\\
$\widehat Z_i$ & Pauli Z operator acting on qubit $i$; a matrix, not a binary variable.\\
$\widehat X_i$ & Pauli X operator acting on qubit $i$; interchanges basis states zero and one.\\
$\Id$ & Identity operator; leaves a vector unchanged.\\
$h_i,J_{ij}$ & Coefficients multiplying Z and ZZ in the Ising cost operator.\\
$C_I$ & Constant identity coefficient after converting to Ising form; called \code{offset}.\\
$\HF$ & Full cost operator $C_I\Id+\HC$; its eigenvalue on a booking basis state is $f(\boldsymbol{x})$.\\
$\HC$ & Nonconstant cost Hamiltonian used to construct gates.\\
$s,\widetilde H_C$ & Positive normalization scale and normalized nonconstant cost $\HC/s$.\\
$H_{\mathrm{Had}}$ & Hadamard gate. Different from the cost Hamiltonian despite the shared letter H.\\
$\ket\psi$, $\bra\psi$ & State vector (ket) and conjugate-transposed row vector (bra).\\
$a_{\boldsymbol{x}}$ & Complex amplitude of outcome $\boldsymbol{x}$; its squared magnitude is a probability.\\
$\langle u\mid v\rangle$ & Inner product of vectors, yielding a complex scalar.\\
$\bra\psi\widehat A\ket\psi$ & Expectation of operator $\widehat A$ in normalized state $\ket\psi$.\\
$\ii$ & Imaginary unit, $\ii^2=-1$; distinct from an italic index $i$.\\
$\phi$, $\theta$ & Phase or rotation angles, measured in radians.\\
$\gamma_\ell$, $\beta_\ell$ & Cost and mixer parameters for layer $\ell$, pronounced gamma and beta.\\
$p$, $\ell$ & Number of QAOA layers and the current layer index.\\
$U_C,U_M$ & Cost and mixer gate operations, each a unitary operator.\\
$\RX,\RZ,\RZZ$ & X rotation, Z rotation, and two-qubit ZZ rotation.\\
$\otimes$ & Tensor product: combines state spaces or operators for multiple qubits.\\
$\sum$, $\prod$ & Add or multiply the indexed quantities. In operator products the rightmost acts first.\\
$\Pr$, $\mathbb E$ & Probability and expectation (probability-weighted mean).\\
Basis state & A definite bit string, such as all qubits zero.\\
Superposition & A state expressed as a linear combination of basis states.\\
Phase & Direction of a complex amplitude; relative phases can affect later interference.\\
Global phase & The same unit-magnitude factor multiplying every amplitude; does not affect measurement probabilities.\\
Interference & Reinforcement or cancellation when operations combine amplitude contributions.\\
Eigenvalue & Number by which an operator multiplies a particular eigenvector.\\
Hamiltonian & A Hermitian operator. Here it represents optimization costs, not battery energy.\\
Ground state & An eigenstate with the smallest eigenvalue of a Hamiltonian.\\
Ansatz & The selected family of parameterized circuits; here repeated cost/mixer layers.\\
Shot & One circuit preparation, execution, and measurement.\\
Estimator & Computes an observable's expectation for a circuit. Our Aer estimator uses ideal precision zero.\\
Statevector & The array of $2^n$ complex amplitudes used by the ideal simulator.\\
COBYLA & Classical derivative-free optimizer that changes circuit angles.\\
\end{longtable}
\renewcommand{\arraystretch}{1}

\clearpage
\section{Mathematical background and gate calculations}\label{app:math}
\subsection{Complex numbers and magnitudes}
A complex number is $a=u+\ii v$ with real $u,v$. Its conjugate is
$a^*=u-\ii v$, and
\begin{equation}
|a|^2=a^*a=u^2+v^2.
\end{equation}
Euler's identity writes a unit-length phase factor as
\begin{equation}
e^{\ii\phi}=\cos\phi+\ii\sin\phi,\qquad |e^{\ii\phi}|=1.
\end{equation}
Thus multiplying an amplitude by $e^{\ii\phi}$ changes its direction but not its
measurement probability. Multiplying \emph{every} amplitude by the same phase
leaves every future measurement probability unchanged under the same subsequent
operations, since the common scalar factors through those linear operations.

For a normalized state $\ket\psi$, $\langle\psi\mid\psi\rangle=1$. Its bra is
the transpose with all entries conjugated. A Hermitian matrix satisfies
$A^\dagger=A$ and has real eigenvalues. A unitary matrix satisfies
$U^\dagger U=\Id$ and preserves normalization. Hamiltonians and gates serve
different roles even though both are operators.

\subsection{Relative phase can be revealed by a later gate}
Define
\begin{equation}
\ket+=\frac{\ket0+\ket1}{\sqrt2},\qquad
\ket-=\frac{\ket0-\ket1}{\sqrt2}.
\end{equation}
Both have immediate measurement probabilities $1/2,1/2$. But
\begin{equation}
H_{\mathrm{Had}}\ket+=\ket0,\qquad
H_{\mathrm{Had}}\ket-=\ket1.
\end{equation}
For the first case, the new zero amplitude is $(1+1)/2=1$ and the new one
amplitude is $(1-1)/2=0$. For the second, those reinforce/cancel roles reverse.
This demonstrates interference without claiming that phase alone changes the
probabilities of an immediate computational-basis measurement.

\subsection{Tensor products and Qiskit's basis order}
Two independent states combine as $\ket a\otimes\ket b$, usually abbreviated
$\ket{ab}$. Their joint vector has four entries. For two qubits, conventional
Qiskit displayed basis order is
\begin{equation}
\ket{q_1q_0}\in\{\ket{00},\ket{01},\ket{10},\ket{11}\}.
\end{equation}
Hence $\widehat Z_0=\Id\otimes\widehat Z$ and
$\widehat Z_1=\widehat Z\otimes\Id$ in that order. For a logical vector
$\boldsymbol{x}=(x_0,\ldots,x_{n-1})$, its statevector array index is
\begin{equation}
k=\sum_{i=0}^{n-1}2^i x_i.
\end{equation}
For $(0,1,0,0,1)$, $k=2+16=18$, displayed in binary as \code{10010}.

\subsection{Matrix forms of the rotations}
For an angle $\theta$, the gates used here are
\begin{align}
\RZ(\theta)&=
\begin{pmatrix}e^{-\ii\theta/2}&0\\0&e^{+\ii\theta/2}\end{pmatrix},\\
\RX(\theta)&=
\begin{pmatrix}
\cos(\theta/2)&-\ii\sin(\theta/2)\\
-\ii\sin(\theta/2)&\cos(\theta/2)
\end{pmatrix},\\
\RZZ(\theta)&=\operatorname{diag}
\left(e^{-\ii\theta/2},e^{+\ii\theta/2},e^{+\ii\theta/2},e^{-\ii\theta/2}\right).
\end{align}
The RZZ diagonal is ordered as $00,01,10,11$. Equal bits have ZZ eigenvalue
$+1$, while different bits have eigenvalue $-1$. The original $x_i x_j$
penalty is not represented by RZZ alone: its converted Z terms and identity
offset also contribute. Together they reproduce a penalty only on $11$.

\subsection{The one-qubit phase/mixer example, fully worked}
After H and $\RZ(-\pi/2)$ the state is
\begin{equation}
\ket{\psi_1}=\frac1{\sqrt2}
\left(e^{\ii\pi/4}\ket0+e^{-\ii\pi/4}\ket1\right).
\end{equation}
Both amplitudes still have magnitude $1/\sqrt2$. The next gate is
\begin{equation}
\RX(-\pi/2)=\frac1{\sqrt2}
\begin{pmatrix}1&\ii\\\ii&1\end{pmatrix}.
\end{equation}
The resulting amplitudes are
\begin{align}
a_0&=\frac{e^{\ii\pi/4}+\ii e^{-\ii\pi/4}}2=e^{\ii\pi/4},\\
a_1&=\frac{\ii e^{\ii\pi/4}+e^{-\ii\pi/4}}2=0.
\end{align}
The final state is $e^{\ii\pi/4}\ket0$, with the same measurement probabilities
as $\ket0$. The first amplitude reinforces while the second cancels. The
remaining global phase is unobservable in the final probabilities.

\begin{lstlisting}[language=Python]
from math import pi
from qiskit import QuantumCircuit
from qiskit.quantum_info import Statevector

qc = QuantumCircuit(1)
qc.h(0)
print(Statevector.from_instruction(qc).probabilities_dict())
qc.rz(-pi / 2, 0)
print(Statevector.from_instruction(qc).probabilities_dict())
qc.rx(-pi / 2, 0)
print(Statevector.from_instruction(qc).probabilities_dict())
\end{lstlisting}
Floating-point output may show a tiny nonzero value such as $10^{-32}$ instead
of exact zero. This is numerical roundoff, not a meaningful probability of failure.

\clearpage
\section{Full numerical audit of the sample objective}\label{app:numeric}
\subsection{Expanded QUBO coefficients}
At $P=57.65$, the generated coefficients are
\begin{align}
C_0&=115.30,\\
(L_0,L_1,L_2,L_3,L_4)&=(-35.60,-31.00,-24.65,-20.65,-38.00),\\
Q_{ij}&=115.30\quad(0\le i<j\le3),\\
Q_{04}&=57.65.
\end{align}
All omitted pairs have coefficient zero. In full,
\begin{align}
f(\boldsymbol{x})={}&115.30-35.60x_0-31x_1-24.65x_2-20.65x_3-38x_4\notag\\
&+115.30(x_0x_1+x_0x_2+x_0x_3+x_1x_2+x_1x_3+x_2x_3)\notag\\
&+57.65x_0x_4.
\end{align}
This expression and Equation~\eqref{eq:sample-qubo} are identical on binary
inputs. The negative coefficients arise from algebraically expanded penalties.

\subsection{Derive one Z coefficient and one pair coefficient}
For qubit zero, incident pair coefficients are $Q_{01},Q_{02},Q_{03},Q_{04}$.
Therefore
\begin{align}
h_0&=-\frac{-35.60}{2}
-\frac{115.30+115.30+115.30+57.65}{4}\notag\\
&=17.80-100.8875=-83.0875.
\end{align}
The A--B conflict yields $J_{04}=Q_{04}/4=14.4125$.
The identity offset is
\begin{align}
C_I&=C_0+\frac12\sum_iL_i+\frac14\sum_{i<j}Q_{ij}\notag\\
&=115.30-74.95+187.3625=227.7125.
\end{align}

\subsection{Verify the selected schedule three ways}
For $\boldsymbol{x}=(0,1,0,0,1)$:
\begin{enumerate}
\item From the factored scheduling expression: all penalties vanish and
      $f=26.65+19.65=46.30$.
\item From the expanded QUBO: only $x_1,x_4$ are one, with no pair between
      them, so $f=115.30-31.00-38.00=46.30$.
\item From Ising signs: $\boldsymbol z=(+1,-1,+1,+1,-1)$. The Z sum is $-167$;
      the six A--A ZZ products sum to zero; the A--B ZZ term is $-14.4125$.
      Thus $E_C=-181.4125$ and $f=C_I+E_C=46.30$.
\end{enumerate}
All representations assign the same score. This is the invariant tested by
checking every one of the 32 basis states.

\subsection{General coefficient generation and a matrix convention trap}
For each vehicle group the compiler adds $P$ to the constant, $-P$ to every
member's linear coefficient, and $2P$ to every within-group pair. Different-car
resource conflicts add $P$ to their pairs. Thus new data changes the coefficients
without a new hand-written equation.

The code stores each pair once, with $i<j$. If instead writing a symmetric
matrix expression $\boldsymbol{x}^{\mathsf T}M\boldsymbol{x}$, remember that
$M_{ij}$ and $M_{ji}$ both contribute. For this convention choose
$M_{ii}=L_i$ and $M_{ij}=M_{ji}=Q_{ij}/2$. Copying upper-triangular coefficients
into both halves without halving would double every pair cost.

\section{Exercises and common points of confusion}
Try these before reading Appendix~\ref{app:answers}.
\begin{enumerate}[label=\textbf{\arabic*.}]
\item A new car C has options $x_5,x_6$. Write its exactly-one penalty. Suppose
      $x_1$ and $x_5$ are different cars using the same station and slot; write
      the additional conflict term.
\item For A, what are the exactly-one penalties when zero, one, two, or three
      options are selected? Why does ``all bits zero'' fail?
\item Evaluate the sample objective at $(1,0,0,0,1)$ using $P=57.65$.
      Which penalties vanish and which does not?
\item A drives an additional $3\,\mathrm{km}$ at $0.2\,\mathrm{kWh/km}$ with
      the same battery target. How much extra charging energy is needed?
\item If $x=1$, what is $z$? Convert $f(x)=8x$ into a Z-form cost operator
      and verify both bit-state scores.
\item Expand $P(a+b-1)^2$ using binary identities. Why is the coefficient of
      $ab$ equal to $2P$, rather than $P$?
\item After a uniform initial state and only diagonal cost gates, what are
      the five-qubit measurement probabilities? Why is the mixer needed?
\item A Z term has coefficient $h=-3$, scale $s=6$, and $\gamma=0.4$.
      What RZ angle should the code use? What RX angle corresponds to $\beta=0.2$?
\item Decode Qiskit's displayed string \code{10100} into logical booking order
      and identify the selected bookings.
\item Does increasing \code{--shots} change the current exact training objective
      or the learned angles when all training settings and the seed stay fixed?
\item Why can a run's best valid sample score be $46.30$ while its expected QUBO
      score is about $73.88$? What does the invalid mode \code{00001} omit?
\item If no valid schedule is sampled, can we conclude that the mathematical
      problem is infeasible? What is different if a car has no individual option?
\end{enumerate}

\subsection{Interpretations to keep separate}
\textbf{A variable is a booking, not a vehicle.} Four of the sample bits belong
to A, but the exactly-one constraint allows only one of those four to be selected.

\textbf{Joint planning is different from simultaneous charging.} The optimizer
selects a complete schedule for all cars. Resource constraints determine which
cars may physically charge at the same time.

\textbf{A penalty is a scoring device.} The implemented mixer can visit invalid
strings. A penalty discourages them but does not make their probability zero.
The final decoder is still necessary.

\textbf{A cost Hamiltonian is not a measurement of electricity.} Its eigenvalues
are mathematical scores. Its constant, signs, and units need not resemble battery
energy or monetary spending.

\textbf{Superposition does not expose all answers at once.} Each measurement
returns one string. QAOA tries to shape their distribution through interference.

\textbf{Correct encoding does not imply successful optimization.} Tests can
prove that an objective matches rules on a small instance while the optimizer
still trains poorly. Evaluate encoding correctness and search quality separately.

\section{Exercise answers}\label{app:answers}
\begin{enumerate}[label=\textbf{\arabic*.}]
\item Add $P(x_5+x_6-1)^2$ for C and $Px_1x_5$ for that resource conflict.
      The builder generates both from the relevant group and pair data.
\item The penalties are $P,0,P,4P$. All-zero bits leave cars unassigned;
      cheap electricity spending alone would otherwise favor serving nobody.
\item Each car has exactly one booking, so its group penalty is zero. Both
      use North early, so the conflict penalty is $P$. The score is
      $22.05+19.65+57.65=99.35$.
\item Extra consumption is $3(0.2)=0.6\,\mathrm{kWh}$. Arrival energy falls by
      that amount, so charging energy increases by $0.6\,\mathrm{kWh}$.
\item $z=-1$. The operator is $4\Id-4\widehat Z$, giving $4-4=0$ on bit zero
      and $4+4=8$ on bit one.
\item $P(a+b-1)^2=P-Pa-Pb+2Pab$. The cross term in the square is $2ab$;
      binary identities reduce $a^2,b^2$ but do not remove that factor of two.
\item Each probability remains $1/32$. Cost gates multiply each amplitude
      by a unit-magnitude phase. Mixing combines contributions with different
      phases and can change the magnitudes, hence probabilities.
\item The RZ angle is $2\gamma h/s=2(0.4)(-3)/6=-0.4$ radians. The RX angle
      is $2\beta=0.4$ radians. The two rotations have different axes and roles.
\item Reverse the displayed string to get $(0,0,1,0,1)$. Select A at South
      early and B at North early, with booking score $33+19.65=52.65$.
\item No. In this implementation the training expectation is ideal and
      independent of final shots. Increasing shots changes the final histogram
      and may reveal previously unobserved outcomes.
\item The expectation averages every outcome, including costly invalid ones.
      The best sample describes one selected valid schedule. \code{00001}
      selects only $x_0$, giving A a booking but none to B.
\item No sampled valid outcome does not prove infeasibility. In contrast,
      an empty option group proves that no complete schedule can serve that
      vehicle under the current model and supplied inputs.
\end{enumerate}

\section{Commands, defaults, and reproducibility}
\begin{center}
\begin{tabular}{lr}
\toprule
CLI setting & Default\\
\midrule
\code{--layers} & 2\\
\code{--iterations} & 160 per restart\\
\code{--restarts} & 3\\
\code{--shots} & 4096\\
\code{--seed} & 7\\
\code{--penalty} & automatic $U+1$\\
\code{--show-circuit} & off\\
\bottomrule
\end{tabular}
\end{center}
The seed fixes initial-angle generation and final simulator sampling for the
current implementation. Optimizer and simulator changes across library versions
can change results even with the same seed. The recorded histogram in this guide
is an observed snapshot, not a permanent expected output specification.

Training requires positive integer layers, shots, evaluation budget, and
restarts; the budget must be at least $2p+2$. The seed must be a nonnegative
integer and a custom penalty must be positive and finite. Larger scenarios
increase qubit requirements; a small number of vehicles does not by itself
guarantee a small quantum circuit if each vehicle has many possible bookings.

For experiments, record a row containing
\begin{equation*}
\begin{aligned}
\text{settings}&=(p,\ \text{budget},\ \text{restarts},\ N,\ \text{seed},\ P),\\
\text{results}&=(\widehat p_{\mathrm{valid}},\ \text{best sample score},\ F_{\mathrm{best}},\ \text{runtime}).
\end{aligned}
\end{equation*}
Keep electricity bill and waiting separately as well. If you change the waiting
weight, the definition of the score changes; scores from different weights are
not directly comparable as if they measured the same preference.

\clearpage
\phantomsection
\addcontentsline{toc}{section}{References}
\begin{thebibliography}{9}
\raggedright
\bibitem{qaoa}
IBM Quantum Documentation, \emph{Quantum approximate optimization algorithm}.
\url{https://quantum.cloud.ibm.com/docs/en/tutorials/quantum-approximate-optimization-algorithm}.
Reference for the QUBO-to-Hamiltonian and hybrid QAOA workflow.

\bibitem{rz}
IBM Quantum Documentation, \emph{RZGate}.
\url{https://quantum.cloud.ibm.com/docs/en/api/qiskit/qiskit.circuit.library.RZGate}.
Reference for the rotation convention $\RZ(\theta)=e^{-\ii\theta Z/2}$.

\bibitem{rzz}
IBM Quantum Documentation, \emph{RZZGate}.
\url{https://quantum.cloud.ibm.com/docs/en/api/qiskit/qiskit.circuit.library.RZZGate}.
Reference for the two-qubit ZZ rotation convention and diagonal form.

\bibitem{aer}
Qiskit Aer Documentation, \emph{EstimatorV2}.
\url{https://qiskit.github.io/qiskit-aer/stubs/qiskit_aer.primitives.EstimatorV2.html}.
Reference for the simulator primitive used during training.

\bibitem{cobyla}
SciPy Documentation, \emph{minimize(method='COBYLA')}.
\url{https://docs.scipy.org/doc/scipy/reference/optimize.minimize-cobyla.html}.
Reference for the optimizer options and evaluation budget.
\end{thebibliography}
